\section{مدیریت ویژگی‌های پیوسته و دسته‌ای}

% Slide 37
\begin{frame}{مدیریت ویژگی‌های پیوسته و دسته‌ای}{\lr{Handling Continuous and Categorical Attributes}}
\begin{itemize}\setlength{\itemsep}{5pt}
    \item چگونه تحلیل وابستگی را بر روی ویژگی‌هایی اعمال کنیم که \textbf{متغیرهای دوگانهٔ نامتقارن نیستند}؟
    \item \textbf{مثال قاعدهٔ ترکیبی:}
    \begin{center}
    \begin{latin}
    \[
    \{\text{Number of Pages} \in [5,10) \wedge (\text{Browser}=\text{Firefox})\} \to \{\text{Buy} = \text{No}\}
    \]
    \end{latin}
    \end{center}
\end{itemize}

\vspace{0.1cm}
\begin{center}
\setlength{\tabcolsep}{3pt}
\fontsize{7.5}{9.5}\selectfont
\begin{tabular}{|c|c|c|c|c|c|c|}
\hline
\rowcolor{mannheimblue!15}
\textbf{شناسه} & \textbf{کشور} & \textbf{مدت جلسه (ثانیه)} & \textbf{صفحات دیده‌شده} & \textbf{جنسیت} & \textbf{مرورگر} & \textbf{خرید} \\ \hline
1 & USA & 982 & 8 & Male & Chrome & No \\ \hline
2 & China & 811 & 10 & Female & Chrome & No \\ \hline
3 & USA & 2125 & 45 & Female & Firefox & Yes \\ \hline
4 & Germany & 596 & 4 & Male & IE & Yes \\ \hline
5 & Australia & 123 & 9 & Male & Firefox & No \\ \hline
\dots & \dots & \dots & \dots & \dots & \dots & \dots \\ \hline
\end{tabular}
\end{center}
\end{frame}

% Slide 38
\begin{frame}{مدیریت ویژگی‌های دسته‌ای}{\lr{Handling Categorical Attributes}}
\begin{itemize}\setlength{\itemsep}{4pt}
    \item \textbf{تبدیل ویژگی‌های دسته‌ای به متغیرهای دودویی نامتقارن:}
    \begin{itemize}\setlength{\itemsep}{2pt}
        \item تعریف یک «آیتم» جدید به‌ازای هر جفت مشخص (ویژگی، مقدار).
        \item برای مثال، جایگزینی ویژگی «نوع مرورگر» با:
        \begin{itemize}\footnotesize
            \item ویژگی: $\text{Browser Type} = \text{Chrome}$
            \item ویژگی: $\text{Browser Type} = \text{Firefox}$
            \item ویژگی: $\text{Browser Type} = \text{Safari}$ و \dots
        \end{itemize}
    \end{itemize}
    \vspace{0.15cm}
    \item \textbf{چالش‌های پیش‌رو:}
    \begin{itemize}\setlength{\itemsep}{3pt}
        \item \alert{\textbf{اگر ویژگی دارای مقادیر ممکن بسیار زیادی باشد:}}
        \begin{itemize}\footnotesize
            \item بسیاری از مقادیر ویژگی ممکن است دارای پشتیبان بسیار پایینی باشند.
            \item \textbf{راهکار احتمالی:} تجمیع مقادیر کم‌پشتیبان در گروه‌های کلی‌تر.
        \end{itemize}
        \item \alert{\textbf{اگر توزیع مقادیر ویژگی دارای چولگی شدید باشد:}}
        \begin{itemize}\footnotesize
            \item برای مثال: \lr{95\%} کاربران دارای $\text{Buy} = \text{No}$ هستند.
            \item اغلب آیتم‌ها با $(\text{Buy}=\text{No})$ همبسته خواهند شد (قواعد بدیهی).
            \item \textbf{راهکار احتمالی:} حذف یا نادیده گرفتن آیتم بسیار پرتکرار.
        \end{itemize}
    \end{itemize}
\end{itemize}
\end{frame}

% Slide 39
\begin{frame}{مدیریت ویژگی‌های پیوسته}{\lr{Handling Continuous Attributes}}
\begin{itemize}\setlength{\itemsep}{4pt}
    \item \textbf{تبدیل متغیرهای پیوسته با گسسته‌سازی:} \lr{(Discretization)}
    \begin{itemize}\setlength{\itemsep}{2pt}
        \item دسته‌بندی با پهنای مساوی \lr{(Equal-width binning)}
        \item دسته‌بندی با بسامد مساوی \lr{(Equal-frequency binning)}
    \end{itemize}
    \vspace{0.15cm}
    \item \textbf{چالش:} طول بازه‌های گسسته‌سازی مستقیماً بر روی پشتیبان و اطمینان اثر می‌گذارد:
    \begin{latin}
    \begin{itemize}\small\setlength{\itemsep}{1pt}
        \item $\{\text{Refund}=\text{No}, (\text{Income}=\$51{,}251)\} \to \{\text{Cheat}=\text{No}\}$
        \item $\{\text{Refund}=\text{No}, (60\text{K} \le \text{Income} \le 80\text{K})\} \to \{\text{Cheat}=\text{No}\}$
        \item $\{\text{Refund}=\text{No}, (0\text{K} \le \text{Income} \le 1\text{B})\} \to \{\text{Cheat}=\text{No}\}$
    \end{itemize}
    \end{latin}
    \vspace{0.15cm}
    \item \textbf{اگر بازه‌ها بیش از حد کوچک باشند:}
    \begin{itemize}\footnotesize
        \item مجموعه‌آیتم‌ها پشتیبان کافی کسب نمی‌کنند و نادیده گرفته می‌شوند.
    \end{itemize}
    \item \textbf{اگر بازه‌ها بیش از حد بزرگ باشند:}
    \begin{itemize}\footnotesize
        \item قواعد اطمینان کافی یا معنای عملی نخواهند داشت.
    \end{itemize}
\end{itemize}
\end{frame}
